\chapter{Zorn, Hartogs i el teorema del bon ordre}\label{cap:zorn}

\begin{objectius}
\apartat{Prerequisits} Del capítol~\ref{cap:bons-ordres}: bons ordres, segments, rigidesa i comparació. Del capítol~\ref{cap:recursio}: inducció i recursió transfinites amb codomini fix, i la comparabilitat total. Del capítol~\ref{cap:relacions}: l'axioma d'elecció, els quocients i l'efectivitat. Del capítol~\ref{cap:families}: seccions, funcions d'elecció i unions indexades.
\apartat{Objectius} En acabar el capítol, el lector hauria de poder:
\begin{itemize}
  \item codificar un bon ordre sobre una part d'\(A\) com un parell \((S,R)\) dins \(P(A)\times P(A\times A)\), i explicar per què el suport ha de ser una dada;
  \item construir \(\Hart(A)\) com a quocient de \(\WO(A)\) per isomorfisme, demostrar que és un bon ordre i que el segment anterior a un tipus té aquell tipus;
  \item demostrar el teorema de Hartogs, sense elecció, i veure'l com l'anàleg ordinal del teorema de Cantor;
  \item demostrar el teorema del bon ordre amb una recursió \emph{acotada per \(\Hart(A)\)}, i explicar per què la cota evita la recursió sobre tots els ordinals;
  \item distingir maximal de màxim, demostrar el lema de Zorn amb la mateixa cota i aplicar-lo (existència de bases);
  \item demostrar el cercle d'equivalències entre elecció, bon ordre i Zorn, marcant quina direcció introdueix l'elecció.
\end{itemize}
\end{objectius}

Els dos capítols anteriors han donat la teoria dels bons ordres i la
recursió transfinita \emph{sense} usar l'elecció. Aquest capítol hi torna a
introduir l'elecció (axioma~\ref{ax:eleccio}) i en treu les tres grans
conseqüències que, en la teoria clàssica, apareixen encadenades: el
\emph{lema de Zorn}, el \emph{teorema del bon ordre} i, com a frontissa
tècnica entre tots dos, el \emph{teorema de Hartogs}. La novetat respecte del
tractament habitual és doble. Primer, el punt de vista estructural manté totes
les construccions \emph{acotades}: en lloc de parlar de \enquote{tots els
ordinals}, treballarem amb parells \((S,R)\), on \(S\) és una part d'un
conjunt fix \(A\) i \(R\) un bon ordre sobre \(S\); aquests parells formen una
part de \(P(A)\times P(A\times A)\), i els seus tipus d'ordre s'obtenen
per un quocient efectiu. Segon, la construcció de Hartogs ---un bon ordre
\(\Hart(A)\) que no s'injecta en \(A\)--- proporciona la \emph{cota} que atura
les recursions: gràcies a ella, el teorema del bon ordre no necessita cap
recursió \enquote{sobre tots els ordinals}, sinó una de sola, acotada per
\(\Hart(A)\). Al final del capítol tancarem el cercle clàssic mostrant que
elecció, teorema del bon ordre i lema de Zorn són equivalents.


\section{El teorema de Hartogs}\label{sec:hartogs}

Comencem per Hartogs, perquè no usa l'elecció i perquè és la peça que fa
acotades les construccions posteriors. La idea és reunir tots els
\emph{tipus} de bons ordres que es poden posar sobre parts d'\(A\), i veure
que aquesta col·lecció, ordenada, és un bon ordre massa gran per cabre dins
\(A\).

\begin{definicio}[Bons ordres sobre parts d'un conjunt fix]\index{Bons ordres sobre parts d'un conjunt fix}\label{def:bons-ordres-parts}
Sigui \(A\) un conjunt. Un \emph{bon ordre sobre una part d'\(A\)} és un
\emph{parell} \((S,R)\) format per un \emph{suport} \(S\hookrightarrow A\)
---un element de \(P(A)\)--- i una relació \(R\hookrightarrow S\times S\)
---un element de \(P(S\times S)\hookrightarrow P(A\times A)\)
(definició~\ref{def:relacio})--- tal que \(R\) és un bon ordre estricte
(definició~\ref{def:bon-ordre}) sobre \emph{tot} \(S\). Codifiquem el parell
com un element del conjunt fix
\[
(S,R)\ \in\ P(A)\times P(A\times A).
\]
La col·lecció d'aquests bons ordres és una part interna
\[
\WO(A)\ \hookrightarrow\ P(A)\times P(A\times A),
\]
tallada per les condicions internes \enquote{\(R\hookrightarrow S\times S\), i
\(R\) és irreflexiva, transitiva, total i ben fundada sobre \(S\)} ---totes
expressables amb igualtats i quantificadors acotats sobre
\(P(A)\times P(A\times A)\) (capítols~\ref{cap:relacions}
i~\ref{cap:parts-interna}). Escrivim \(\camp(S,R):=S\) per al suport, que ara
és una dada del parell i no una quantitat derivada.
\end{definicio}

\begin{observacio}[Per què el parell, i no la sola relació]\label{obs:parell-suport}
Codificar el bon ordre només per la relació \(R\) i recuperar-ne el suport
com \(\camp(R)=\{a: \exists b\,(a\mathrel R b\vee b\mathrel R a)\}\) falla en
els casos degenerats. Sobre un singletó \(\{a\}\), un bon ordre estricte és la
relació \emph{buida}, i aleshores \(\camp(R)=\varnothing\), de manera que la
relació buida sola no distingeix l'ordre buit sobre \(\varnothing\) dels
ordres sobre subconjunts d'un element. El parell \((S,R)\) fixa el suport
explícitament i elimina l'ambigüitat: el bon ordre buit sobre \(\{a\}\) queda
representat per \((\{a\},\varnothing)\), diferent de
\((\varnothing,\varnothing)\). Amb la codificació antiga, en canvi,
\(\Hart(A)\) d'un singletó \(A\) podria contenir un únic tipus que
\emph{sí} s'injecta en \(A\), contradient l'enunciat de Hartogs
(teorema~\ref{teo:hartogs}); el parell ho evita.
\end{observacio}

\begin{observacio}[Tot és acotat]\label{obs:hartogs-acotat}
Aquesta és la clau del tractament estructural de Hartogs. Un bon ordre sobre
una part d'\(A\) \emph{no} és un objecte extern: és un element del conjunt fix
\(P(A)\times P(A\times A)\), i \(\WO(A)\) n'és una part. No parlem de
\enquote{tots els bons ordres} en abstracte, sinó dels que caben sobre parts
d'un conjunt donat, que formen un conjunt ja construït. Tota la resta de la
secció ---la identificació per isomorfisme, el quocient, l'ordre--- passa dins
aquest ambient, sense cap classe global.
\end{observacio}

\begin{definicio}[Tipus d'ordre; el conjunt de Hartogs]\index{Tipus d'ordre; el conjunt de Hartogs}\label{def:hartogs}
Sobre \(\WO(A)\), la relació \enquote{\((S,R)\) i \((S',R')\) són isomorfs com
a bons ordres} ---és a dir, hi ha una bijecció \(S\to S'\) que transporta
\(R\) en \(R'\)--- és una relació d'equivalència interna \(\cong\) (reflexiva,
simètrica i transitiva, per composició d'isomorfismes d'ordre). Per
l'efectivitat de les relacions d'equivalència (axioma~\ref{ax:exactitud})
existeix el quocient
\[
\Hart(A)\ :=\ \WO(A)/\cong,
\]
el \emph{conjunt dels tipus d'ordre} dels bons ordres sobre parts d'\(A\). La
classe \([(S,R)]\in\Hart(A)\) és el \emph{tipus d'ordre} del bon ordre
\((S,R)\); quan el suport se sobreentén escrivim simplement \([R]\).
Ordenem \(\Hart(A)\) per la relació de segment inicial: per a dos tipus
\(\alpha=[(S,R)]\) i \(\beta=[(S',R')]\),
\[
\alpha<\beta
\quad:\Longleftrightarrow\quad
(S,R)\ \text{és isomorf a un segment inicial propi de}\ (S',R')
\]
(definició~\ref{def:mes-curt}), que no depèn dels representants:
component isomorfismes d'ordre, si \((S,R)\cong(S_1,R_1)\) i
\((S',R')\cong(S'_1,R'_1)\), un isomorfisme de \((S,R)\) sobre un segment
propi de \((S',R')\) dona un isomorfisme de \((S_1,R_1)\) sobre un segment
propi de \((S'_1,R'_1)\).
\end{definicio}

\begin{proposicio}[\(\Hart(A)\) és un bon ordre]\label{prop:hartogs-bon-ordre}
La relació \(<\) sobre \(\Hart(A)\) de la definició~\ref{def:hartogs} és un
bon ordre.\par
\emph{Axiomes usats:} conjunt de parts, relacions i quocients efectius
(capítols~\ref{cap:relacions} i~\ref{cap:parts-interna},
axioma~\ref{ax:exactitud}), comparabilitat de bons ordres
(teorema~\ref{teo:comparabilitat}); \emph{cap elecció}.
\end{proposicio}

\begin{prova}
Que \(<\) és un ordre estricte total és la comparació de bons ordres del
capítol~\ref{cap:recursio}: dos tipus qualssevol són comparables
(teorema~\ref{teo:comparabilitat}), i l'exclusivitat i la transitivitat són
la tricotomia i la proposició~\ref{prop:comparacio-ordres}; la
irreflexivitat és que cap bon ordre no és isomorf a un segment propi seu
(corol·lari~\ref{cor:no-segment-propi}). Falta el bon ordre: tota part no
inicial \(T\hookrightarrow\Hart(A)\) té mínim. Sigui
\(\alpha_0=[(S_0,R_0)]\in T\) qualsevol. Si \(\alpha_0\) ja és el mínim, fet.
Si no, els tipus de \(T\) estrictament menors que \(\alpha_0\) corresponen,
per definició de \(<\), a segments inicials propis de \((S_0,R_0)\); i els
segments inicials de \((S_0,R_0)\) estan en bijecció creixent amb els punts
del suport \(S_0\) (proposició~\ref{prop:parts-inicials}), que és un bon
ordre. Per tant \(\{\beta\in T:\beta<\alpha_0\}\) es correspon amb una part no
inicial de \(S_0\), que té mínim pel bon ordre \(R_0\); aquest mínim dona el
tipus mínim de \(T\). Així \(\Hart(A)\) és ben
ordenat.
\end{prova}

\begin{lema}[El segment anterior a un tipus té aquell tipus]\label{lema:segment-tipus}
Sigui \(\alpha=[(S,R)]\in\Hart(A)\). El segment inicial propi
\[
\Hart(A)_{<\alpha}\ :=\ \{\,\beta\in\Hart(A):\beta<\alpha\,\}
\]
amb l'ordre \(<\) heretat és un bon ordre de tipus exactament \(\alpha\); és a
dir, \((\Hart(A)_{<\alpha},<)\) és isomorf a \((S,R)\).\par
\emph{Axiomes usats:} els de la proposició~\ref{prop:hartogs-bon-ordre};
\emph{cap elecció}.
\end{lema}

\begin{prova}
Per definició de \(<\), un tipus \(\beta=[(S',R')]\) compleix \(\beta<\alpha\)
si i només si \((S',R')\) és isomorf a un segment inicial propi de \((S,R)\).
Ara bé, els segments inicials propis de \((S,R)\) estan en bijecció creixent
amb els punts de \(S\): a cada \(s\in S\) li correspon el segment
\(S_{<s}=\{t\in S: t\mathrel R s\}\), i tot segment inicial propi és d'aquesta
forma (proposició~\ref{prop:parts-inicials}). Dos punts diferents donen
segments no isomorfs, perquè cap bon ordre no és isomorf a un segment propi
seu (corol·lari~\ref{cor:no-segment-propi}). Així l'aplicació
\[
s\ \longmapsto\ [\,(S_{<s},\,R\!\restriction\!S_{<s})\,]
\]
és una bijecció de \(S\) sobre \(\Hart(A)_{<\alpha}\), i és creixent perquè
\(s\mathrel R s'\) equival a \(S_{<s}\) ser segment inicial propi de
\(S_{<s'}\). Per tant és un isomorfisme d'ordre \((S,R)\cong
\Hart(A)_{<\alpha}\), i \(\Hart(A)_{<\alpha}\) té tipus \(\alpha\).
\end{prova}

\begin{teorema}[Hartogs]\index{Hartogs, teorema de}\label{teo:hartogs}
Per a tot conjunt \(A\), el bon ordre \(\Hart(A)\) \emph{no} admet cap
monomorfisme \(\Hart(A)\hookrightarrow A\). És a dir, \(\Hart(A)\not\preceq A\):
hi ha un bon ordre estrictament \enquote{massa llarg} per a \(A\).\par
\emph{Axiomes usats:} els de la proposició~\ref{prop:hartogs-bon-ordre}; cap
axioma nou; \emph{cap elecció}.
\end{teorema}

\begin{prova}
Suposem, cap a una contradicció, un monomorfisme \(m\colon\Hart(A)\hookrightarrow
A\). Transportant l'ordre de \(\Hart(A)\) per \(m\), la imatge
\(S_{*}:=m(\Hart(A))\hookrightarrow A\) és una part d'\(A\) dotada d'un bon
ordre \(R_{*}\hookrightarrow S_{*}\times S_{*}\) isomorf a \(\Hart(A)\) ---és a
dir, \((S_{*},R_{*})\in\WO(A)\), i el seu tipus
\([(S_{*},R_{*})]\in\Hart(A)\) és el tipus del bon ordre \(\Hart(A)\) mateix.
Ara bé, per construcció, \(\Hart(A)\) ordena \emph{tots} els tipus de bons
ordres sobre parts d'\(A\), inclòs \([(S_{*},R_{*})]\); i pel
lema~\ref{lema:segment-tipus} el segment inicial de \(\Hart(A)\) determinat per
un tipus \(\alpha\) és isomorf a un bon ordre de tipus \(\alpha\). En
particular, \(\Hart(A)\) seria isomorf al segment inicial de \(\Hart(A)\)
determinat per \([(S_{*},R_{*})]\), un segment inicial \emph{propi} de si
mateix ---perquè \([(S_{*},R_{*})]\) és un element de \(\Hart(A)\). Això
contradiu que cap bon ordre no és isomorf a cap segment propi seu
(corol·lari~\ref{cor:no-segment-propi}). Per tant no existeix cap tal \(m\), i
\(\Hart(A)\not\preceq A\).
\end{prova}

\begin{observacio}[Hartogs com a paral·lel de Cantor]\label{obs:hartogs-cantor}
El teorema de Hartogs és, per als \emph{ordres}, l'anàleg del que el teorema
de Cantor (teorema~\ref{teo:cantor}) és per a les \emph{mides}: així com
\(P(A)\) produeix una mida estrictament més gran que \(A\), \(\Hart(A)\)
produeix un bon ordre estrictament més llarg que qualsevol que càpiga dins
\(A\). Cap dels dos no usa elecció, i tots dos són \enquote{diagonals} en un
sentit ampli: la contradicció surt de fer actuar la construcció sobre si
mateixa. La diferència de funció és que \(P\) mesura cardinalitat i
\(\Hart\) mesura longitud d'ordre; per això \(\Hart(A)\) serà la cota natural
de les recursions que segueixen, on el que importa és la longitud.
\end{observacio}

\section{El teorema del bon ordre}\label{sec:teorema-bon-ordre}

Amb l'elecció i la cota de Hartogs podem provar que tot conjunt admet un bon
ordre. La idea és anar exhaurint \(A\) triant successivament un element del
que queda; el que garanteix que el procés \emph{acaba} ---i per tant que és
una recursió acotada, no una \enquote{sobre tots els ordinals}--- és
\(\Hart(A)\).

\begin{teorema}[Teorema del bon ordre (Zermelo)]\label{teo:bon-ordre}
Suposem l'axioma d'elecció (axioma~\ref{ax:eleccio}). Aleshores tot conjunt
\(A\) admet un bon ordre: existeix una relació \(R\hookrightarrow A\times A\)
que és un bon ordre sobre tot \(A\).\par
\emph{Axiomes usats:} elecció (axioma~\ref{ax:eleccio}), Hartogs
(teorema~\ref{teo:hartogs}), recursió transfinita amb codomini fix
(teorema~\ref{teo:recursio-transfinita}); cap recursió \enquote{sobre tots els
ordinals} ---la construcció és acotada per \(\Hart(A)\).
\end{teorema}

\begin{prova}
Per l'elecció (axioma~\ref{ax:eleccio}), la família de les parts no inicials
d'\(A\) admet una \emph{funció d'elecció}: una aplicació
\(c\colon P(A)\setminus\{\varnothing\}\to A\) amb \(c(S)\in S\) per a tota
part no inicial \(S\) ---una secció de la família corresponent
(definició~\ref{def:seccio-familia}, proposició~\ref{prop:epi-regular-eleccio}).
Afegim un punt exterior \(\top\), formant \(A^{+}=A+\mathbf1\), i definim per
recursió transfinita (teorema~\ref{teo:recursio-transfinita}) sobre el bon
ordre \(\Hart(A)\) una aplicació \(F\colon\Hart(A)\to A^{+}\) amb la regla:
donada la història \(F|_{\Hart(A)_{<\alpha}}\),
\[
F(\alpha)\ :=\
\begin{cases}
c\bigl(A\setminus F(\Hart(A)_{<\alpha})\bigr) & \text{si } A\setminus F(\Hart(A)_{<\alpha})\neq\varnothing,\\[1mm]
\top & \text{altrament,}
\end{cases}
\]
és a dir, a cada etapa es tria ---amb \(c\)--- un element d'\(A\) encara no
fet servir, i \(\top\) quan ja no en queda cap. El codomini \(A^{+}\) és
\emph{fix}, de manera que el teorema~\ref{teo:recursio-transfinita} s'aplica i
\(F\) existeix.

\emph{\(F\) s'atura.} Sobre la part \(H_0=\{\alpha:F(\alpha)\neq\top\}\), la
regla fa \(F\) \emph{injectiva} cap a \(A\) (encara no hi ha cap ordre fixat
sobre \(A\), de manera que no té sentit dir-la \enquote{creixent}): si
\(\beta<\alpha\) i \(F(\alpha)\neq\top\), aleshores
\(F(\alpha)\notin F(\Hart(A)_{<\alpha})\), en particular \(F(\alpha)\neq
F(\beta)\). Si \(F\) no prengués mai \(\top\),
\(F\colon\Hart(A)\hookrightarrow A\) seria un monomorfisme, contra el teorema
de Hartogs (teorema~\ref{teo:hartogs}). Per tant existeix algun \(\alpha\) amb
\(F(\alpha)=\top\); sigui \(\alpha_0\) el mínim
(bon ordre de \(\Hart(A)\), proposició~\ref{prop:hartogs-bon-ordre}).

\emph{El bon ordre sobre \(A\).} Per la minimalitat de \(\alpha_0\),
\(F(\alpha)\neq\top\) per a tot \(\alpha<\alpha_0\); i per a
\(\alpha>\alpha_0\) també \(F(\alpha)=\top\), perquè
\(F(\Hart(A)_{<\alpha})\supseteq F(\Hart(A)_{<\alpha_0})=A\) i el complement
és buit. Així \(F(\alpha)=\top\) exactament per a \(\alpha\ge\alpha_0\); i
\(F(\alpha_0)=\top\) vol dir \(A\setminus F(\Hart(A)_{<\alpha_0})=\varnothing\),
és a dir \(F(\Hart(A)_{<\alpha_0})=A\). Així la restricció
\[
F|_{\Hart(A)_{<\alpha_0}}\colon\ \Hart(A)_{<\alpha_0}\ \longrightarrow\ A
\]
és injectiva i exhaustiva: una bijecció entre el segment inicial
\(\Hart(A)_{<\alpha_0}\) ---un bon ordre (proposició~\ref{prop:segment-ben-ordenat})---
i \(A\) (encara sense ordre). Transportant l'ordre del segment per aquesta
bijecció s'obté un bon ordre \(R\hookrightarrow A\times A\) sobre tot \(A\), i
respecte d'aquest ordre transportat la bijecció esdevé, ja sí, un isomorfisme
d'ordre creixent. El procés ha estat
\emph{acotat} per \(\Hart(A)\): no hem recorregut \enquote{tots els
ordinals}, sinó el bon ordre concret \(\Hart(A)\), i ens hem aturat en
\(\alpha_0<\) la seva longitud.
\end{prova}

\begin{observacio}[Per què cal la cota de Hartogs]\label{obs:cota-hartogs}
La construcció \enquote{anar triant un element del que queda} és la prova
clàssica de Zermelo, però tal com se sol enunciar amaga una recursió sobre
una col·lecció d'ordinals que no és, a priori, un conjunt. El tractament
estructural ho evita fent la recursió sobre el bon ordre \emph{concret}
\(\Hart(A)\), que sí que és un conjunt (definició~\ref{def:hartogs}): la
recursió transfinita amb codomini fix (teorema~\ref{teo:recursio-transfinita})
s'hi aplica directament, i el teorema de Hartogs garanteix que abans
d'esgotar \(\Hart(A)\) ja s'han esgotat els elements d'\(A\). Aquesta és la
raó tècnica per la qual Hartogs precedeix el teorema del bon ordre: és la cota
que converteix una recursió \enquote{potencialment sense fi} en una recursió
acotada de Nivell~I més elecció.
\end{observacio}

\section{El lema de Zorn}\label{sec:zorn}

La tercera forma de maximalitat és el lema de Zorn, que en la pràctica
matemàtica és la més usada. El demostrarem a partir del teorema del bon
ordre; primer, el vocabulari.

\begin{definicio}[Cadena, cota superior, maximal]\index{Cadena, cota superior, maximal}\label{def:zorn-nocions}
Sigui \((P,\le)\) un conjunt parcialment ordenat (un ordre reflexiu,
antisimètric i transitiu, capítol~\ref{cap:bons-ordres}).
\begin{itemize}
\item Una \emph{cadena} de \(P\) és una part \(C\hookrightarrow P\) totalment
ordenada per \(\le\): dos elements qualssevol de \(C\) són comparables.
\item Una \emph{cota superior} d'una part \(S\hookrightarrow P\) és un
element \(b\in P\) amb \(s\le b\) per a tot \(s\in S\).
\item Un element \(m\in P\) és \emph{maximal} si cap element no el supera
estrictament: \(m\le x\) implica \(m=x\).
\end{itemize}
\end{definicio}

\begin{observacio}[Maximal no és màxim]\label{obs:maximal-maxim}
Convé no confondre \emph{maximal} amb \emph{màxim}. Un \emph{màxim} és un
element \(\ge\) que \emph{tots} els altres; un \emph{maximal} és només un
element que \emph{cap} altre no supera. En un ordre total coincideixen, però
en un ordre parcial pot haver-hi diversos maximals i cap màxim: per exemple,
en les parts pròpies d'un conjunt de dos elements ordenades per inclusió, els
dos singletons són tots dos maximals, i no hi ha màxim. El lema de Zorn dona
un \emph{maximal}, no un màxim; afirmar que dona un màxim seria fals.
\end{observacio}

\begin{teorema}[Lema de Zorn]\index{Zorn, lema de}\label{teo:zorn}
Suposem l'axioma d'elecció (axioma~\ref{ax:eleccio}). Sigui \((P,\le)\) un
conjunt parcialment ordenat no inicial en què \emph{tota cadena té una cota
superior}. Aleshores \(P\) té almenys un element maximal.\par
\emph{Axiomes usats:} elecció (axioma~\ref{ax:eleccio}), teorema del bon
ordre (teorema~\ref{teo:bon-ordre}) i recursió transfinita
(teorema~\ref{teo:recursio-transfinita}); tot acotat, cap recursió sobre una
classe.
\end{teorema}

\begin{prova}
Suposem, cap a una contradicció, que \(P\) no té cap maximal: per a tot
\(x\in P\) existeix \(x'\) amb \(x<x'\). Aleshores tota cadena \(C\) té, a
més d'una cota superior, una cota superior \emph{estricta} ---un element
estrictament més gran que tots els de \(C\)---: prenent una cota superior
\(b\) de \(C\) i un \(b'>b\) (que existeix perquè \(b\) no és maximal),
\(b'\) supera estrictament tot \(C\).

Considerem la família indexada per les parts de \(P\) que té, sobre cada
part \(C\), la part de \(P\) de les seves cotes superiors estrictes
(una part de \(P(P)\times P\), tallada amb quantificadors interns). Les
fibres sobre les cadenes són no buides pel paràgraf anterior; restringint
la família a la part de les cadenes i aplicant l'elecció
(axioma~\ref{ax:eleccio}) obtenim una aplicació \(c\) que a cada cadena
\(C\) assigna una cota superior estricta \(c(C)\notin C\). Fixem també un
punt qualsevol \(p_0\in P\) (\(P\) no és inicial). Definim per recursió
transfinita (teorema~\ref{teo:recursio-transfinita}) sobre \(\Hart(P)\) una
aplicació \(F\colon\Hart(P)\to P\) amb la regla \emph{total}
\[
F(\alpha):=
\begin{cases}
c\bigl(F(\Hart(P)_{<\alpha})\bigr) & \text{si } F(\Hart(P)_{<\alpha})\text{ és una cadena},\\[1mm]
p_0 & \text{altrament},
\end{cases}
\]
on la condició \enquote{ser una cadena} és interna (quantificadors acotats
sobre \(P\times P\)); la segona clàusula no s'arribarà a usar mai, però fa
que la regla estigui definida per a tota història, com el
teorema~\ref{teo:recursio-transfinita} exigeix. Per inducció transfinita
(teorema~\ref{teo:induccio-transfinita}) es comprova alhora que la imatge
\(F(\Hart(P)_{<\alpha})\) és sempre una cadena i que \(F\) és estrictament
creixent: si ho és per a tot \(\beta<\alpha\), els valors \(F(\beta)\)
estan totalment ordenats, de manera que la imatge és una cadena, el primer
cas de la regla s'aplica, i \(F(\alpha)=c(\cdot)\) supera estrictament tots
els \(F(\beta)\) amb \(\beta<\alpha\). Per tant \(F\colon\Hart(P)\to P\) és un
monomorfisme (una aplicació estrictament creixent des d'un ordre total és
injectiva), és a dir \(\Hart(P)\preceq P\). Però això contradiu el teorema de
Hartogs (teorema~\ref{teo:hartogs}), \(\Hart(P)\not\preceq P\). La
contradicció ve de suposar que \(P\) no té maximal; per tant \(P\) en té un.
\end{prova}

\begin{observacio}[La cota de Hartogs, altre cop]\label{obs:zorn-hartogs}
La prova de Zorn és, estructuralment, la mateixa que la del teorema del bon
ordre: una recursió que tria a cada pas un element \enquote{nou} (aquí, una
cota estricta) i que \emph{no pot continuar indefinidament} perquè
\(\Hart(P)\) no s'injecta en \(P\). És el teorema de Hartogs, i no cap
recursió sobre una classe d'ordinals, el que atura el procés i dona la
contradicció. Aquesta és la manera estructural, i acotada, de fer el que en
la prova clàssica es fa \enquote{iterant fins que ja no es pot}.
\end{observacio}

\subsection{Una aplicació: tot espai vectorial té una base}\label{sub:base}

Perquè Zorn no quedi com un enunciat purament fundacional, en donem una
aplicació matemàtica genuïna: l'existència de bases. Es formula amb el
vocabulari d'aquest llibre prenent un espai vectorial com un conjunt amb les
operacions internes corresponents; el que importa aquí és l'estructura de
l'argument de Zorn.

\begin{proposicio}[Existència de base]\label{prop:base}
Suposem l'elecció. Tot espai vectorial \(V\) sobre un cos \(K\) té una base:
una part \(B\hookrightarrow V\) linealment independent i generadora.\par
\emph{Axiomes usats:} lema de Zorn (teorema~\ref{teo:zorn}); per tant elecció.
\end{proposicio}

\begin{prova}
Sigui \(P\) el conjunt de les parts linealment independents de \(V\),
ordenat per inclusió ---una part de \(P(V)\) tallada per la condició interna
d'independència lineal (cap combinació lineal finita no trivial no s'anul·la),
que és un conjunt parcialment ordenat no inicial (la part buida hi és, o
qualsevol singletó no nul). Vegem que tota cadena \(C\hookrightarrow P\) té
cota superior: la unió \(U=\bigcup C\) (unió indexada de parts,
capítol~\ref{cap:families}) és linealment independent, perquè tota
combinació lineal involucra només un nombre \emph{finit} de vectors, que per
ser \(C\) una cadena viuen tots en un mateix membre de la cadena, on són
independents; i \(U\) conté tots els membres de \(C\), és a dir n'és cota
superior. Pel lema de Zorn (teorema~\ref{teo:zorn}), \(P\) té un element
maximal \(B\). Aquesta \(B\) és una base: és independent per ser a \(P\), i
és generadora perquè si algun \(v\in V\) no fos combinació lineal de \(B\),
aleshores \(B\cup\{v\}\) seria independent i estrictament més gran, contra la
maximalitat de \(B\).
\end{prova}

\begin{observacio}[Altres aplicacions]\label{obs:zorn-aplicacions}
El mateix esquema ---ordenar per inclusió, veure que les unions de cadenes
resten dins la classe, i prendre un maximal--- dona molts resultats
d'existència no constructius: l'extensió d'un filtre a un ultrafiltre, la
clausura algebraica d'un cos, l'existència d'ideals maximals en un anell amb
unitat, o el teorema de Hahn--Banach en la seva forma algebraica. En tots,
Zorn substitueix una tria \enquote{fins a l'infinit} per un maximal; i en
tots, la propietat de tancament per unions de cadenes és el que cal
comprovar.
\end{observacio}

\section{Les tres equivalències}\label{sec:equivalencies}

Tanquem el capítol amb el cercle clàssic: sobre la base dels altres axiomes
del llibre, les tres afirmacions són equivalents.

\begin{teorema}[Elecció, bon ordre i Zorn són equivalents]\label{teo:equivalencies}
Sobre els axiomes del llibre \emph{sense} l'elecció, les tres afirmacions
següents són equivalents:
\[
\begin{aligned}
\text{(E) elecció}
\quad&\Longleftrightarrow\quad
\text{(BO) tot conjunt admet un bon ordre}\\
&\Longleftrightarrow\quad
\text{(Z) lema de Zorn.}
\end{aligned}
\]
\emph{Axiomes usats:} recursió transfinita i Hartogs
(capítols~\ref{cap:recursio} i aquest), que no depenen de l'elecció; les
implicacions que \emph{introdueixen} l'elecció es marquen a la prova.
\end{teorema}

\begin{prova}
\emph{(E) \(\Rightarrow\) (BO)} és el teorema del bon ordre
(teorema~\ref{teo:bon-ordre}); \emph{(E) \(\Rightarrow\) (Z)} és el lema de
Zorn (teorema~\ref{teo:zorn}). Falten dues direccions que \emph{retornen}
l'elecció.

\emph{(BO) \(\Rightarrow\) (E).} Sigui \(p\colon E\twoheadrightarrow I\) una
família de parts no inicials (un epimorfisme regular). Per (BO), fixem un bon
ordre sobre \(E\). Aleshores cada fibra \(E_i=p^{-1}(i)\) és una part no
inicial d'\(E\), i té un \emph{mínim} per al bon ordre
(definició~\ref{def:bon-ordre}); l'aplicació \(s\colon I\to E\) que a cada
\(i\) assigna \(\min(E_i)\) és una secció de \(p\) ---està ben definida
perquè el mínim és únic i canònic, sense cap tria addicional. Per tant tota
família de parts no buides té secció, que és l'elecció (axioma~\ref{ax:eleccio}).

\emph{(Z) \(\Rightarrow\) (BO)}, i per tant \((Z)\Rightarrow(E)\). Sigui \(A\)
un conjunt; volem ben ordenar-lo. Considerem el conjunt \(P\) dels
\emph{bons ordres parcials} d'\(A\): els parells \((S,R)\) amb
\(S\hookrightarrow A\) i \(R\) un bon ordre sobre \(S\) ---una part de
\(P(A)\times P(A\times A)\). Ordenem \(P\) per \emph{extensió inicial}:
\((S,R)\le(S',R')\) si \(S\hookrightarrow S'\), \(R\) és la restricció de
\(R'\) a \(S\), i \(S\) és un \emph{segment inicial} de \((S',R')\). Tota cadena \(C\) té cota superior: siguin
\(S_{C}=\bigcup\{S:(S,R)\in C\}\) i \(R_{C}=\bigcup\{R:(S,R)\in C\}\) les
unions indexades (capítol~\ref{cap:families}). \(R_{C}\) és un ordre
estricte total sobre \(S_{C}\), perquè dos punts qualssevol són en un mateix
membre de la cadena (que està totalment ordenada per extensió inicial) i
allà són comparables; i és un bon ordre: si \(T\hookrightarrow S_{C}\) és
una part no buida, prenem un punt \(t\in T\) i un membre \((S,R)\in C\) amb
\(t\in S\); el mínim de \(T\cap S\) per \(R\) és el mínim de tot \(T\),
perquè qualsevol altre punt de \(T\) o bé és a \(S\), o bé és en un membre
que estén \((S,R)\) \emph{inicialment}, i aleshores és \(\ge\) tots els
punts de \(S\). Finalment \((S_{C},R_{C})\) estén inicialment cada membre
de \(C\), de manera que n'és cota superior. (La cadena buida té per cota
\((\varnothing,\varnothing)\in P\).) Pel lema de Zorn
(teorema~\ref{teo:zorn}), \(P\) té un maximal \((S_{*},R_{*})\). Ha de ser
\(S_{*}=A\): si quedés un element \(a\in A\setminus S_{*}\), afegint-lo al
final de \(R_{*}\) obtindríem un bon ordre parcial estrictament més gran,
contra la maximalitat. Per tant \(R_{*}\) és un bon ordre sobre tot \(A\), i
val (BO).

Les tres afirmacions són, doncs, equivalents. Cap direcció no ha exigit una
recursió més forta que la del capítol~\ref{cap:recursio}: la construcció de
\(\Hart\) i les recursions són totes acotades i de codomini fix.
\end{prova}

\begin{resumcap}
\apartat{Idees centrals}
\begin{enumerate}
  \item Un bon ordre sobre una part d'\(A\) és un parell \((S,R)\) dins \(P(A)\times P(A\times A)\); el suport ha de ser una dada perquè la relació buida no el determina. \(\WO(A)\) és una part interna i \(\Hart(A)=\WO(A)/\!\cong\) existeix per l'efectivitat.
  \item \(\Hart(A)\), ordenat per segment inicial, és un bon ordre, i el segment anterior a un tipus \(\alpha\) té tipus \(\alpha\). Teorema de Hartogs: \(\Hart(A)\not\preceq A\), sense elecció; és l'anàleg ordinal de Cantor.
  \item Teorema del bon ordre: amb una funció d'elecció sobre les parts no buides d'\(A\) i una recursió \emph{acotada per \(\Hart(A)\)}, s'exhaureix \(A\) abans d'esgotar \(\Hart(A)\), i la bijecció resultant transporta el bon ordre del segment.
  \item Lema de Zorn amb la mateixa cota: si no hi hagués maximal, una recursió que tria cotes estrictes donaria \(\Hart(P)\preceq P\), contra Hartogs. Zorn dona un maximal, no un màxim.
  \item Elecció, bon ordre i Zorn són equivalents; (BO)\(\Rightarrow\)(E) pel mínim de cada fibra, (Z)\(\Rightarrow\)(BO) per maximalitat en els bons ordres parcials.
\end{enumerate}
\apartat{Errors típics} Codificar un bon ordre només per la relació, perdent el suport; parlar de \enquote{tots els ordinals} en lloc de \(\Hart(A)\); donar una regla de recursió definida només en el cas bo (ha de ser total); confondre maximal amb màxim, o cota superior amb cota estricta; i oblidar que la unió d'una cadena de bons ordres parcials és ben ordenada \emph{perquè} les extensions són inicials.
\apartat{Lectura recomanada} \cite[§§15--17]{halmos} per a l'elecció, Zorn i el teorema del bon ordre en la teoria clàssica; \cite[cap.~5]{jech2003} per al teorema de Hartogs, l'aleph de Hartogs i les equivalències de l'elecció; \cite[cap.~I.12]{kunen2011} per a la presentació axiomàtica d'aquests resultats; i \cite[apèndix~B]{lawvere-rosebrugh} per a la formulació estructural de l'elecció i els seus equivalents.
\end{resumcap}
